\chapter{Background}
\label{ch:background}

This chapter gives the background needed for the rest of the report and reviews existing applications and research. It concludes in the gap that MedLift-3D fills.

\section{Preliminaries}

This section presents the physics of X-ray projection and tomosynthesis, iterative reconstruction and null space, diffusion models, Gaussian representations and nodule mass, conformal prediction, and explainability and detection measures which are used throughout the remainder of the report.

\subsection{X-ray Projection and Digital Tomosynthesis}
 A detector pixel records the line integral of the linear attenuation coefficient $\mu$ along the ray from the source. With $I_0$ incident photons the expected count is
\begin{equation}
  \bar I = I_0 \exp\Big(-\int_{\text{ray}} \mu(\mathbf{r})\, \mathrm{d}l\Big), \qquad I \sim \mathrm{Poisson}(\bar I),
\end{equation}
so fewer photons (lower dose) mean more noise. CT values are expressed in Hounsfield units, $\mathrm{HU} = 1000\,(\mu-\mu_{\text{water}})/\mu_{\text{water}}$.

\paragraph{The missing wedge.} CT collects projections over \ang{180} or more. Chest DTS moves the tube along a short linear track, typically $\pm15^\circ$ (GE VolumeRAD) to $\pm20^\circ$ (Shimadzu SonialVision), with 15--74 views~\cite{dobbins2003dts}. By the Fourier slice theorem each view measures one plane of the object's 3D Fourier transform, so a narrow sweep leaves a wedge of depth frequencies unmeasured (\cref{fig:dts}b). Reconstructions are sharp in the coronal plane but every structure is smeared in depth.

\subsection{Iterative Reconstruction and the Null Space}
 Writing the projection operator as $A$, the measured line integrals as $\mathbf{b}$ and the volume as $\mathbf{x}$, SIRT~\cite{gilbert1972sirt,herman2009} solves $A\mathbf{x}\approx\mathbf{b}$ by repeated normalised back-projection of the residual:
\begin{equation}
  \mathbf{x}^{(k+1)} = \mathbf{x}^{(k)} + \lambda\, C A^{\top} R \big(\mathbf{b} - A\mathbf{x}^{(k)}\big),
\end{equation}
where $R$ and $C$ hold the inverse row and column sums of $A$. It is stable under noise but cannot recover anything that lies in the null space of $A$.

\paragraph{Seen and unseen parts.} Any volume splits uniquely as $\mathbf{x} = \mathbf{x}_{\text{seen}} + \mathbf{x}_{\text{unseen}}$ with $\mathbf{x}_{\text{seen}} \in \mathcal{R}(A^{\top})$ and $\mathbf{x}_{\text{unseen}} \in \mathcal{N}(A)$, i.e.\ $A\mathbf{x}_{\text{unseen}}=0$. The X-rays determine only the seen part; the unseen part of any reconstruction comes from its prior~\cite{bhadra2021,wang2023ddnm}. \Cref{fig:seen} illustrates this.

\begin{figure}[htbp]
  \centering
  \resizebox{0.82\linewidth}{!}{% Seen (range) / unseen (null-space) decomposition and what each reconstruction contributes
\begin{tikzpicture}[font=\footnotesize, >=Stealth]
  \draw[->, thick, edge] (0,0) -- (8.2,0) node[below=4pt, anchor=north east, align=right] {seen component $\mathbf{x}_{\mathrm{seen}} \in \mathcal{R}(A^{\!\top})$\\ fixed by the X-rays};
  \draw[->, thick, edge] (0,0) -- (0,5.2) node[above, align=center] {unseen component $\mathbf{x}_{\mathrm{unseen}} \in \mathcal{N}(A)$\\ (missing wedge)};
  % data-consistent set
  \draw[accent, very thick] (5,-0.3) -- (5,4.9);
  \node[accent, align=left, anchor=west] at (5.15,4.7) {all volumes with $A\mathbf{x} = \mathbf{b}$\\ (agree with the 15 views)};
  % SIRT / min norm
  \fill[edge] (5,0) circle (0.11);
  \node[anchor=south west, align=left] at (5.15,0.12) {SIRT / minimum norm:\\ no unseen content};
  % truth
  \node[star, star points=5, fill=red!70!black, inner sep=1.6pt] (gt) at (5,3.1) {};
  \node[anchor=east, text=red!70!black] at (4.75,3.1) {true CT $\mathbf{x}^\ast$};
  % prior samples
  \foreach \y in {2.35,2.6,2.8,2.95,3.25,3.45,3.7,3.95} { \fill[accent!70] (5,\y) circle (0.075); }
  \draw[decorate, decoration={brace, amplitude=5pt}] (5.35,3.98) -- (5.35,2.32) node[midway, right=7pt, align=left] {8 diffusion samples:\\ unseen content from the prior};
  % skill annotation
  \draw[<->, densely dashed, edge] (4.3,0.05) -- (4.3,3.05);
  \node[anchor=east, align=right, text=quiet] at (4.2,1.55) {unseen skill: does the\\ prior's unseen part\\ correlate with the truth's?};
\end{tikzpicture}
}
  \caption[Seen and unseen components of a reconstruction]{Seen (range) and unseen (null-space) components. All volumes on the vertical line agree with the measured views. SIRT returns the minimum-norm point with no unseen content; diffusion samples supply unseen content, which we test against the true CT (unseen skill, \cref{sec:method}).}
  \label{fig:seen}
\end{figure}

\subsection{Denoising Diffusion Models}
 A diffusion model learns to reverse a gradual noising process~\cite{ho2020ddpm}. With a cosine noise schedule $\bar\alpha_t$~\cite{nichol2021improved}, a noisy sample is
\begin{equation}
  \mathbf{x}_t = \sqrt{\bar\alpha_t}\,\mathbf{x}_0 + \sqrt{1-\bar\alpha_t}\,\boldsymbol\epsilon, \qquad \boldsymbol\epsilon\sim\mathcal{N}(\mathbf{0},I).
\end{equation}
We train the network to predict the velocity $\mathbf{v} = \sqrt{\bar\alpha_t}\,\boldsymbol\epsilon - \sqrt{1-\bar\alpha_t}\,\mathbf{x}_0$ (v-prediction~\cite{salimans2022v}), minimising
\begin{equation}
  \mathcal{L}(\theta) = \mathbb{E}_{\mathbf{x}_0,\mathbf{c},t,\boldsymbol\epsilon}\, \big\| \mathbf{v} - \hat{\mathbf{v}}_\theta(\mathbf{x}_t, \mathbf{c}, t)\big\|_2^2 ,
\end{equation}
where the condition $\mathbf{c}$ is the DTS reconstruction, so the model learns $p(\mathbf{x}_0\mid\mathbf{c})$. DDIM~\cite{song2021ddim} samples in tens of steps instead of a thousand, and data-consistency (DC) sampling corrects every denoised estimate with the physics so that samples agree with the measurements~\cite{chung2023dps,chung2024dds}.

\subsection{3D Gaussian Representations and Nodule Mass}
 A volume can be written as a sum of anisotropic Gaussians with centres $\boldsymbol\mu_i$, covariances $\Sigma_i = R_i S_i S_i^{\top} R_i^{\top}$ and amplitudes $a_i$~\cite{kerbl2023gs,zha2024r2}:
\begin{equation}
  f(\mathbf{r}) = \sum_i a_i \exp\!\Big(-\tfrac12(\mathbf{r}-\boldsymbol\mu_i)^{\top}\Sigma_i^{-1}(\mathbf{r}-\boldsymbol\mu_i)\Big),
  \qquad
  M = \int f = \sum_i a_i (2\pi)^{3/2} s_{i1}s_{i2}s_{i3},
  \label{eq:gaussmass}
\end{equation}
where $s_{i1},s_{i2},s_{i3}$ are the Gaussian's standard deviations. The mass is therefore closed form, and the field can be evaluated on any grid.

\paragraph{Nodule mass.} In this report mass is the nodule's excess attenuation over the surrounding lung, converted to milligrams with the CT density scale (about 0.7--0.85~mg/mm$^3$ on LIDC nodules). Mass doubling time is computed like VDT but from mass; it detects growth of part-solid nodules earlier~\cite{dehoop2010}.

\subsection{Conformal Prediction}
 Split conformal prediction~\cite{vovk2022alrw,angelopoulos2023} turns any predictor into intervals with a finite-sample guarantee: if calibration and test cases are exchangeable, a $(1-\alpha)$ interval contains the truth with probability at least $1-\alpha$. A \emph{normalised} version scales each interval by a difficulty estimate (here the spread of the diffusion hypotheses); a \emph{Mondrian} version calibrates within groups (here nodule size groups) so that coverage holds in every group; and \emph{conformal predictive systems} (CPS) return a full calibrated cumulative distribution for each case~\cite{vovk2019cps,bostrom2022crepes}, from which calibrated probabilities such as $P(V\geq 300~\text{mm}^3)$ follow. \Cref{fig:conformal} shows the procedure.

\begin{figure}[htbp]
  \centering
  \resizebox{\linewidth}{!}{% Split-conformal calibration on held-out patients and its use on a new nodule
\begin{tikzpicture}[font=\footnotesize, >=Stealth]
  % (a) calibration scores
  \begin{scope}
    \node[anchor=north west, font=\small\bfseries] at (-0.2,3.55) {(a)};
    \node[anchor=north west, text width=6.8cm] at (0.45,3.55) {calibration scores on 117 held-out patients:\\ $s_i = |\log(V_i/\hat V_i)|\,/\,\hat\sigma_i$};
    \draw[->, edge] (0,0) -- (7.2,0) node[right] {$s$};
    \foreach \x/\h in {0.3/1.2,0.7/1.9,1.1/2.5,1.5/2.3,1.9/1.8,2.3/1.4,2.7/1.0,3.1/0.75,3.5/0.55,3.9/0.4,4.3/0.3,4.7/0.22,5.1/0.15,5.5/0.1,5.9/0.07} {
      \fill[accent!35] (\x-0.18,0) rectangle (\x+0.18,\h);
    }
    \draw[red!70!black, very thick] (4.05,-0.1) -- (4.05,2.4);
    \node[red!70!black, align=left, anchor=west] at (4.15,2.0) {$\hat q$: the $\lceil (n{+}1)(1{-}\alpha)\rceil$-th\\ smallest score};
    \node[text=quiet] at (2.0,-0.35) {90\% of scores};
    \draw[<->, quiet] (0.1,-0.6) -- (4.0,-0.6);
  \end{scope}
  % (b) test interval
  \begin{scope}[xshift=8.6cm]
    \node[anchor=north west, font=\small\bfseries] at (-0.2,3.55) {(b)};
    \node[anchor=north west, text width=6.5cm] at (0.45,3.55) {new nodule: prediction $\hat V$,\\ ensemble spread $\hat\sigma$};
    \draw[->, edge] (0,1.2) -- (6.8,1.2) node[right] {$\log V$};
    \fill[accent!20] (1.4,0.9) rectangle (5.2,1.5);
    \draw[accent, very thick] (1.4,0.85) -- (1.4,1.55) (5.2,0.85) -- (5.2,1.55);
    \fill[accent] (3.3,1.2) circle (0.09);
    \node[below=5pt] at (3.3,1.05) {$\hat V$};
    \node[above=3pt] at (1.4,1.55) {$\ell = \hat V e^{-\hat q\hat\sigma}$};
    \node[above=3pt] at (5.2,1.55) {$u = \hat V e^{+\hat q\hat\sigma}$};
    \node[align=center, text=quiet] at (3.3,-0.35) {$P(V \in [\ell,u]) \geq 1-\alpha = 0.9$\\ for exchangeable patients};
  \end{scope}
  \draw[->, thick, edge] (7.3,1.6) -- (8.4,1.6);
\end{tikzpicture}
}
  \caption[Split-conformal calibration]{Split-conformal calibration (schematic). (a) Normalised scores on held-out calibration patients and their $(1-\alpha)$ quantile $\hat q$. (b) The interval for a new nodule, valid for exchangeable patients without any assumption on the error distribution.}
  \label{fig:conformal}
\end{figure}

\subsection{Explainability and Detection Metrics}
 Input$\times$gradient saliency~\cite{shrikumar2017} multiplies each input voxel by the gradient of the output with respect to it, showing which voxels drove a detector's score. Detection is measured by the Competition Performance Metric (CPM), the mean sensitivity at 1/8, 1/4, 1/2, 1, 2, 4 and 8 false positives per scan~\cite{setio2017luna}.

\section{Literature Review}

The review has two parts: applications that already address parts of the problem in practice, and research on radiation and screening, nodule measurement, chest DTS, deep priors for sparse-view reconstruction, Gaussian tomography, and trustworthiness and uncertainty.

\subsection{Similar Applications}

Parts of the problem are already available in clinical practice, but the tools available do not measure nodules in 3D from DTS. The GE and Shimadzu chest DTS units generate coronal slices which radiologists read visually, and nodule size is either estimated visually or a diameter is provided in a single slice. CT nodule software, offered by scanner vendors and AI companies, segments nodules and report volumes and doubling time -- but only on CT. Detection systems like the ones compared in the LUNA16 challenge predict the locations of nodules but do not measure them. Manual segmentation is possible with general tools, like 3D Slicer, ITK-SNAP, etc., which do not carry out any learnt prior or uncertainty, while the research codes for sparse-view reconstruction generate the whole volumes for CT geometry and report only the image quality. These applications are summarised in \cref{tab:similar} and each has a drawback for our problem.

\begin{table}[!htbp]
\centering\small
\caption{Applications similar to MedLift-3D.}
\label{tab:similar}
\begin{tabularx}{\linewidth}{>{\raggedright\arraybackslash}p{3.6cm} >{\raggedright\arraybackslash}X >{\raggedright\arraybackslash}X}
\toprule
Application & What it does & What it lacks for our problem \\
\midrule
Chest DTS units (GE VolumeRAD, Shimadzu SonialVision) & Acquire and reconstruct DTS slices for reading & No 3D nodule measurement; depth smear left to the reader \\
CT nodule volumetry software (vendor lung packages, AI nodule tools) & Detect and segment nodules on CT, report volume and VDT & Requires CT at every follow-up; no DTS support \\
CT computer-aided detection (LUNA16-style)~\cite{setio2017luna} & Suggest nodule locations on CT & Detection only; no measurement or uncertainty \\
3D Slicer, ITK-SNAP & Viewing and manual segmentation & No learned prior, no calibrated numbers \\
Research sparse-view reconstruction codes (R$^2$-Gaussian, 3DGR-CT)~\cite{zha2024r2,li2025gr} & Reconstruct whole volumes from few views & CT geometry, image metrics only, no clinical numbers \\
\bottomrule
\end{tabularx}
\end{table}

\subsection{Related Research}

\paragraph{Radiation and screening.} Mascalchi and Sali~\cite{mascalchi2017} reviewed the cumulative dose of low-dose CT screening (9.3--13.0~mSv over ten years) and argued for ALARA and ultra-low-dose methods; Rampinelli et al.~\cite{rampinelli2017} estimated the cancer risk from screening and follow-up CT.

\paragraph{3D measurement of nodules.} Lee~\cite{lee2025volumetry} and de Margerie-Mellon et al.~\cite{demargerie2018} showed that volumetry is more reproducible than diameters but varies between software, and Revel et al.~\cite{revel2004twod} showed that 2D measurements of small nodules are unreliable. NELSON~\cite{horeweg2014nelson} and the European position statement~\cite{oudkerk2017eups} define the volume and VDT categories we use. De Hoop et al.~\cite{dehoop2010} showed that mass detects growth of ground-glass nodules earlier than volume or diameter.

\paragraph{Chest DTS.} Vikgren et al.~\cite{vikgren2008} found that DTS detected about three times more nodules than chest radiography, and Dobbins and Godfrey~\cite{dobbins2003dts} described the depth-resolution limit set by the sweep angle. DTS studies measure nodules in 2D or by reader judgement.

\paragraph{Sparse-view CT with deep priors.} Score-based and diffusion models have been applied in the sinogram domain~\cite{guan2024sino}, with wavelet decompositions~\cite{xu2024sword,wu2024wavelet}, for spectral CT~\cite{liu2024spectral} and zero-shot with variational solvers~\cite{he2025vss}; DPS and DDS add data consistency to sampling~\cite{chung2023dps,chung2024dds}, and foundation models learn lung CT priors from \num{100000} scans~\cite{wang2025lctfound}. Classical iterative reconstruction~\cite{herman2009} and compressed sensing~\cite{candes2006} suppress noise but over-smooth fine detail, and implicit neural fields~\cite{mildenhall2020nerf} are slow per patient. Lighter physics-conditioned diffusion~\cite{vavekanand2025}, deep image prior with ADMM~\cite{chang2025dip}, mixed priors for metal artifacts~\cite{jiang2025mixed}, learned view sampling~\cite{yang2025sampling} and unsupervised low-dose denoising~\cite{jeon2023mmnet} address other parts of the problem. All of these report image quality on CT volumes; none produces a calibrated clinical measurement or uses DTS geometry.

\paragraph{3D Gaussians in tomography.} 3DGR-CT~\cite{li2025gr} and R$^2$-Gaussian~\cite{zha2024r2} reconstruct sparse-view CT with Gaussians, and spatiotemporal Gaussians handle motion~\cite{xie2025pmf,liu2025fmgs}. They use Gaussians for speed and quality, not for closed-form mass, and have no learned prior for the missing wedge.

\paragraph{Trustworthiness and uncertainty.} Bhadra et al.~\cite{bhadra2021} defined hallucinations through the null space of the imaging operator, DDNM~\cite{wang2023ddnm} uses a range--null-space decomposition inside diffusion sampling, and Zhang et al.~\cite{zhang2025pail} used uncertainty estimates to suppress artifacts in limited-data CT. Conformal predictive systems are available in the open-source crepes library~\cite{bostrom2022crepes} but have not been applied to imaging biomarkers from DTS.

\begin{table}[!htbp]
\centering\small
\caption{Summary of the most related research.}
\label{tab:related}
\begin{tabularx}{\linewidth}{>{\raggedright\arraybackslash}p{2.6cm} >{\raggedright\arraybackslash}p{2.6cm} >{\raggedright\arraybackslash}X >{\raggedright\arraybackslash}X}
\toprule
Work & Method & Key finding & Gap for DTS nodule measurement \\
\midrule
Li et al.~\cite{li2025gr} & FBP-guided 3D Gaussians & Faster than neural fields & Whole-volume CT; no prior; no measurement \\
Zha et al.~\cite{zha2024r2} & Radiative Gaussian splatting & Accurate volume integration & No learned prior; no uncertainty \\
Guan et al.~\cite{guan2024sino} & Sinogram-domain diffusion & Handles varying sparsity & 2D; no 3D consistency \\
He et al.~\cite{he2025vss} & Zero-shot variational score solver & No paired data needed & Heavy compute; image metrics only \\
Chung et al.~\cite{chung2023dps,chung2024dds} & DPS / DDS sampling & Measurement-consistent samples & Not for DTS; not calibrated \\
Wang and Guo~\cite{wang2025lctfound} & Lung CT foundation model & Strong transfer & CT only \\
Bhadra et al.~\cite{bhadra2021} & Null-space hallucination maps & Formal definition of hallucination & No truth-referenced skill per nodule \\
Zhang et al.~\cite{zhang2025pail} & Progressive artifact learning & Fewer false structures & Uncertainty not calibrated \\
de Hoop et al.~\cite{dehoop2010} & Nodule mass on CT & Mass detects growth earlier & CT only \\
Bostr\"om~\cite{bostrom2022crepes} & Conformal predictive systems & Valid CDFs for any regressor & Not applied to imaging biomarkers \\
\bottomrule
\end{tabularx}
\end{table}

\section{Gap Analysis}
\label{sec:gap}

The review shows four gaps. First, DTS research stops at 2D detection and size; no work argues from the projection physics which 3D quantity DTS can measure. Second, generative sparse-view methods target whole CT volumes and image metrics, not a clinical number for one nodule, and none is built for the DTS sweep. Third, Gaussian tomography uses Gaussians for speed, not for closed-form mass. Fourth, reconstruction uncertainty is usually raw variance, which is not calibrated, and explanations of what a prior invented are not checked against the truth. \Cref{tab:gap} compares the closest works.

\begin{table}[!htbp]
\centering\small
\caption{Domain comparison: \cmark\ addressed, \xmark\ not addressed, $\circ$ partially.}
\label{tab:gap}
\begin{adjustbox}{max width=\linewidth}
\begin{tabular}{l ccccccc}
\toprule
Work & \makecell{DTS\\geometry} & \makecell{Learned\\prior} & \makecell{Data\\consistency} & \makecell{Gaussian\\mass} & \makecell{Clinical\\biomarker} & \makecell{Calibrated\\uncertainty} & \makecell{Seen/unseen\\explanation} \\
\midrule
Li et al.~\cite{li2025gr} & \xmark & \xmark & \cmark & \xmark & \xmark & \xmark & \xmark \\
Zha et al.~\cite{zha2024r2} & \xmark & \xmark & \cmark & \xmark & \xmark & \xmark & \xmark \\
He et al.~\cite{he2025vss} & \xmark & \cmark & \cmark & \xmark & \xmark & \xmark & \xmark \\
Chung et al.~\cite{chung2023dps} & \xmark & \cmark & \cmark & \xmark & \xmark & \xmark & \xmark \\
Zhang et al.~\cite{zhang2025pail} & \xmark & \cmark & \cmark & \xmark & \xmark & $\circ$ & \xmark \\
Bhadra et al.~\cite{bhadra2021} & \xmark & \xmark & -- & \xmark & \xmark & \xmark & \cmark \\
de Hoop et al.~\cite{dehoop2010} & \xmark & \xmark & -- & \xmark & \cmark & \xmark & \xmark \\
\midrule
\textbf{MedLift-3D} & \cmark & \cmark & \cmark & \cmark & \cmark & \cmark & \cmark \\
\bottomrule
\end{tabular}
\end{adjustbox}
\end{table}

\paragraph{What we do not claim.} 3D Gaussians in medical imaging, uncertainty for Gaussian reconstructions, conformal prediction in CT and for lung nodules, null-space hallucination analysis and saliency maps all exist already. Our contribution is their combination for DTS nodule measurement, the mass-first argument, and a calibrated, truth-checked evaluation.

\section{Summary}

DTS is low-dose and affordable but blind in depth, and existing AI reconstruction neither respects which quantities DTS can measure nor states how far its numbers can be trusted. MedLift-3D addresses this gap by measuring mass first, filling depth with a data-consistent diffusion prior, calibrating every number with conformal methods, and explaining what was measured and what was supplied by the prior. \Cref{ch:design} turns this into requirements and a design.
